\documentclass[9pt,a4paper]{extarticle}
\usepackage{parskip}

% =====================
% Codificação, idioma e layout
% =====================
\usepackage[T1]{fontenc}
\usepackage[utf8]{inputenc}
\usepackage[brazil]{babel}
\usepackage{lmodern}
\usepackage{geometry}
\geometry{margin=2cm}
\usepackage{microtype}

% =====================
% Matemática e símbolos
% =====================
\usepackage{amsmath, amssymb, amsthm}
\usepackage{mathtools}
\usepackage{bm}

% =====================
% Links
% =====================
\usepackage{hyperref}
\hypersetup{colorlinks=true,linkcolor=black,citecolor=black,urlcolor=blue}

% =====================
% Outros
% =====================
\usepackage{enumitem}

\begin{document}

\begin{center}
    \large
    Universidade Federal de Ouro Preto (UFOP)\\
    Instituto de Ciências Exatas e Biológicas (ICEB)\\
    Departamento de Computação (DECOM)\\[0.5em]
    \normalsize
    Disciplina: Otimização Não Linear\\[0.5em]
    \textbf{Roteiro de Aula – Métodos de Busca Direta}
\end{center}

\vspace{1em}

\section*{Objetivos de Aprendizagem}

Ao final desta aula, o(a) estudante deverá ser capaz de:
\begin{itemize}
    \item Diferenciar conceitualmente métodos diretos e indiretos de otimização.
    \item Descrever a ideia básica de métodos aleatórios de busca direta:
    \begin{itemize}
        \item Random Jumping;
        \item Random Walk;
        \item Random Walk with Direction Exploitation;
        \item Grid Search.
    \end{itemize}
    \item Aplicar e analisar, em exemplos numéricos simples:
    \begin{itemize}
        \item o método univariado (\emph{univariate search});
        \item o método de Powell (pattern search / direções conjugadas);
        \item o método do simplex de Nelder--Mead.
    \end{itemize}
\end{itemize}

\subsection*{Leitura Recomendada}

\begin{itemize}
    \item Kalyanmoy Deb (ou Rao): capítulo sobre \emph{Direct Search Methods} (seção de métodos aleatórios, método univariado, Powell e simplex).
\end{itemize}

\section{Introdução: Métodos Diretos vs. Indiretos}

\subsection{Motivação Geral}

Em problemas de otimização não linear, podemos agrupar os algoritmos em duas grandes famílias:

\begin{description}
    \item[Métodos indiretos (baseados em derivadas)] utilizam explicitamente informações de derivadas da função objetivo:
    \begin{itemize}
        \item gradiente $\nabla f(\bm{x})$;
        \item Hessiana $\nabla^2 f(\bm{x})$ ou aproximações.
    \end{itemize}
    Exemplos: método do gradiente, Newton, quasi-Newton, gradiente conjugado.
    \item[Métodos diretos (sem derivadas)] utilizam apenas valores da função $f(\bm{x})$, sem exigir derivadas analíticas ou numéricas.
    Exemplos: métodos aleatórios de busca, método univariado, Powell, simplex de Nelder--Mead.
\end{description}

Contextos em que métodos diretos são particularmente úteis:
\begin{itemize}
    \item Funções objetivo \emph{caixa-preta}, oriundas de simulações complexas;
    \item Funções não diferenciáveis, com descontinuidades ou ruído numérico;
    \item Situações em que derivadas são caras ou difíceis de obter e/ou implementar.
\end{itemize}

Contraste importante para enfatizar em sala:
\begin{itemize}
    \item Métodos indiretos tendem a ser mais eficientes quando derivadas são confiáveis;
    \item Métodos diretos tendem a ser mais \textbf{robustos} e mais gerais, porém normalmente exigem mais avaliações de $f$.
\end{itemize}

\section{Visão Geral: Métodos Aleatórios e Grid Search}

Nesta seção, a ideia é apenas dar a \emph{intuição} dos métodos, sem entrar em detalhes de implementação.

\subsection{Random Jumping Method}

\subsubsection*{Ideia Básica}

\begin{itemize}
    \item Suponha um problema irrestrito, mas com limites estabelecidos para cada variável:
    \[
        l_i \le x_i \le u_i,\quad i = 1,2,\dots,n.
    \]
    \item Geramos vetores aleatórios $\bm{r} = (r_1,\dots,r_n)^T$ com $r_i \sim U(0,1)$.
    \item Construímos pontos no hipercubo:
    \[
        x_i = l_i + r_i (u_i - l_i), \quad i = 1,\dots,n.
    \]
    \item Avaliamos $f(\bm{x})$ em cada ponto e mantemos o melhor valor encontrado.
\end{itemize}

Comentários:
\begin{itemize}
    \item Método extremamente simples e fácil de programar;
    \item Adequado como \textbf{exploração inicial} do espaço de busca;
    \item Desvantagem: pode precisar de muitas amostras aleatórias para localizar bons mínimos em alta dimensão.
\end{itemize}

\subsection{Random Walk Method}

\subsubsection*{Ideia Básica}

\begin{itemize}
    \item Em vez de escolher pontos independentes, construímos uma trajetória:
    \[
        \bm{X}_{i+1} = \bm{X}_i + \lambda \bm{u}_i,
    \]
    onde:
    \begin{itemize}
        \item $\lambda$ é um comprimento de passo (step length);
        \item $\bm{u}_i$ é um vetor unitário aleatório.
    \end{itemize}
    \item Se $f(\bm{X}_{i+1}) < f(\bm{X}_i)$, aceitamos o novo ponto;
    \item Caso contrário, geramos uma nova direção aleatória.
    \item Quando não há melhorias após muitas tentativas, reduzimos $\lambda$ (por exemplo, $\lambda \leftarrow \lambda/2$).
\end{itemize}

Observações:
\begin{itemize}
    \item Usa-se um critério geométrico para evitar viés para as diagonais do hipercubo (aceita-se apenas vetores aleatórios com comprimento $R \le 1$ antes da normalização);
    \item O método “caminha” pelo espaço, explorando regiões de menor valor de $f$.
\end{itemize}

\subsection{Random Walk com Exploração de Direção}

\subsubsection*{Motivação}

Se uma direção aleatória $\bm{u}_i$ reduziu o valor de $f$, faz sentido aproveitar ao máximo essa direção, em vez de dar apenas um passo de comprimento fixo.

\subsubsection*{Ideia}

\begin{itemize}
    \item A partir de $\bm{X}_i$, fazemos uma minimização unidimensional ao longo de $\bm{u}_i$:
    \[
        \bm{X}_{i+1} = \bm{X}_i + \lambda_i^\ast \bm{u}_i,
    \]
    com
    \[
        \lambda_i^\ast = \arg\min_\lambda f(\bm{X}_i + \lambda \bm{u}_i).
    \]
    \item Para encontrar $\lambda_i^\ast$, emprega-se um dos métodos unidimensionais vistos anteriormente (seção áurea, interpolação quadrática etc.).
\end{itemize}

\subsection{Grid Search}

\subsubsection*{Ideia Básica}

\begin{itemize}
    \item Discretizamos cada intervalo $[l_i,u_i]$ em $p_i$ pontos:
    \[
        x_i^{(1)}, x_i^{(2)}, \dots, x_i^{(p_i)},\quad i=1,\dots,n.
    \]
    \item Construímos a grade cartesiana com $p_1 p_2 \cdots p_n$ pontos;
    \item Avaliamos $f$ em todos os pontos da grade;
    \item Escolhemos o ponto de menor valor de $f$ como solução aproximada.
\end{itemize}

Comentários:
\begin{itemize}
    \item Simples e determinístico;
    \item Sofre com a explosão combinatória em alta dimensão (ex.: $n=10$, $p_i=4$ leva a $4^{10}$ pontos);
    \item Útil para:
    \begin{itemize}
        \item problemas de baixa dimensão;
        \item geração de bons pontos iniciais para métodos mais sofisticados.
    \end{itemize}
\end{itemize}

\section{Método Univariado (Univariate Search)}

\subsection{Ideia Geral}

No método univariado, otimizamos uma variável de cada vez, mantendo as demais fixas.

\begin{itemize}
    \item Em cada iteração, escolhemos uma direção coordenada $S_i$ (vetor da base canônica);
    \item Fazemos uma busca unidimensional ao longo dessa direção;
    \item Ao final de $n$ buscas (uma por coordenada), completa-se um ciclo.
\end{itemize}

As direções típicas são:
\[
S_1 = (1, 0, 0, \dots, 0)^T,\quad
S_2 = (0, 1, 0, \dots, 0)^T,\quad
\dots,\quad
S_n = (0, 0, \dots, 1)^T.
\]

\subsection{Algoritmo Esquemático}

\begin{enumerate}
    \item Escolher um ponto inicial $\bm{X}_1$ e um \emph{probe} pequeno $\varepsilon$ (por exemplo, $\varepsilon = 0{,}01$).
    \item Para cada direção $S_i$:
    \begin{enumerate}[label=(\alph*)]
        \item Calcular $f(\bm{X}_k)$, $f(\bm{X}_k + \varepsilon S_i)$ e $f(\bm{X}_k - \varepsilon S_i)$;
        \item Determinar a direção de descida:
        \begin{itemize}
            \item se $f(\bm{X}_k + \varepsilon S_i) < f(\bm{X}_k)$, usar $+S_i$;
            \item se $f(\bm{X}_k - \varepsilon S_i) < f(\bm{X}_k)$, usar $-S_i$;
            \item se nenhuma reduz $f$, considerar que não há melhora naquela coordenada.
        \end{itemize}
        \item Se houver direção de descida, resolver
        \[
            \min_\lambda f(\bm{X}_k \pm \lambda S_i)
        \]
        com um método unidimensional adequado e obter $\lambda_i^\ast$.
        \item Atualizar:
        \[
            \bm{X}_{k+1} = \bm{X}_k \pm \lambda_i^\ast S_i.
        \]
    \end{enumerate}
    \item Repetir ciclos sobre as $n$ direções até que a melhoria em $f$ seja pequena.
\end{enumerate}

\subsection{Exemplo Numérico}

Considere a função
\[
    f(x_1,x_2) = x_1 - x_2 + 2x_1^2 + 2x_1x_2 + x_2^2,
\]
e o ponto inicial
\[
    \bm{X}_1 = (0,0)^T.
\]

\subsubsection*{Iteração 1: direção $S_1 = (1,0)^T$}

\begin{itemize}
    \item $f(0,0) = 0$;
    \item $f(\varepsilon,0) = \varepsilon + 2\varepsilon^2$; para $\varepsilon = 0{,}01$,
    \[
        f(0{,}01,0) \approx 0{,}0102;
    \]
    \item $f(-\varepsilon,0) = -\varepsilon + 2\varepsilon^2$;
    \[
        f(-0{,}01,0) \approx -0{,}0098.
    \]
\end{itemize}

Como $f(-\varepsilon,0) < f(0,0)$, a direção de descida é $-S_1$.

Agora minimizamos $f(-\lambda,0) = 2\lambda^2 - \lambda$. Derivando:
\[
    \frac{d}{d\lambda} (2\lambda^2 - \lambda) = 4\lambda - 1 = 0
    \quad\Rightarrow\quad
    \lambda^\ast = \frac{1}{4}.
\]

Logo, o novo ponto é
\[
    \bm{X}_2 = \bm{X}_1 - \lambda^\ast S_1 = \left(-\frac{1}{4}, 0\right)^T,
\]
com
\[
    f\left(-\frac{1}{4},0\right) = -\frac{1}{8}.
\]

\subsubsection*{Iteração 2: direção $S_2 = (0,1)^T$}

A partir de $\bm{X}_2 = (-1/4,0)$, repete-se o procedimento:
\begin{itemize}
    \item calcula-se $f(-1/4,0)$, $f(-1/4,\varepsilon)$ e $f(-1/4,-\varepsilon)$;
    \item determina-se se a direção correta é $+S_2$ ou $-S_2$;
    \item realiza-se a minimização unidimensional $f(-1/4, \lambda)$ para encontrar $\lambda^\ast$;
    \item obtém-se o novo ponto $\bm{X}_3$.
\end{itemize}

\textbf{Comentário didático:} é interessante realizar pelo menos um ciclo completo no quadro, mostrando a trajetória aproximada e discutindo:
\begin{itemize}
    \item lentidão perto do ótimo;
    \item possibilidade de travamento em vales estreitos (exemplo com falha do univariate em um vale muito “fino”).
\end{itemize}

\section{Método de Powell (Pattern Search / Direções Conjugadas)}

\subsection{Intuição}

Problema do método univariado:
\begin{itemize}
    \item As direções de busca são sempre paralelas aos eixos de coordenadas;
    \item Perto do mínimo, isso gera uma trajetória “dente de serra” e convergência lenta;
    \item Em funções com vales alongados, o método pode oscilar muito.
\end{itemize}

Ideia de Powell:
\begin{itemize}
    \item Observar que, para funções quadráticas em duas variáveis, as retas que ligam pontos alternados do método univariado tendem a passar pelo mínimo;
    \item Usar essas retas como novas direções de busca, chamadas \textbf{direções padrão} ou \textbf{direções conjugadas};
    \item Em funções quadráticas de $n$ variáveis, um conjunto de $n$ direções \emph{conjugadas} permite atingir o mínimo em no máximo $n$ passos de minimização unidimensional (convergência quadrática).
\end{itemize}

\subsection{Ideia de Direções Conjugadas}

Seja $A$ uma matriz simétrica $n \times n$ (por exemplo, a Hessiana de uma função quadrática). Um conjunto de vetores $S_1,\dots,S_n$ é \textbf{$A$-conjugado} se
\[
    S_i^T A S_j = 0,\quad \text{para } i \neq j.
\]

Direções ortogonais (no sentido usual) são um caso particular, quando $A = I$.

\subsection{Esboço do Algoritmo de Powell}

Para um problema em $n$ variáveis:

\begin{enumerate}
    \item Escolher ponto inicial $\bm{X}_1$.
    \item Inicializar as direções $S_1,\dots,S_n$ como os vetores coordenados.
    \item \textbf{Para cada ciclo}:
    \begin{enumerate}[label=(\alph*)]
        \item Guardar o ponto inicial do ciclo: $Z = X$.
        \item Para $i=1,\dots,n$:
        \begin{itemize}
            \item Minimizar $f$ ao longo de $S_i$ a partir de $X$;
            \item Atualizar o ponto $X$ com o resultado da minimização unidimensional.
        \end{itemize}
        \item Ao final do ciclo, construir a \textbf{direção padrão}:
        \[
            S_p = X - Z.
        \]
        \item Minimizar $f$ ao longo de $S_p$, atualizando novamente $X$.
        \item Atualizar o conjunto de direções substituindo a \emph{mais antiga} por $S_p$.
    \end{enumerate}
    \item Repetir ciclos até que a melhoria em $f$ seja pequena.
\end{enumerate}

\subsection{Exemplo Numérico (Esboço)}

Usando novamente
\[
    f(x_1,x_2) = x_1 - x_2 + 2x_1^2 + 2x_1x_2 + x_2^2,
\]
e ponto inicial $\bm{X}_1 = (0,0)$:

\begin{itemize}
    \item Primeiramente, execute um ciclo de busca univariada ao longo de $S_1$ e $S_2$ para obter um ponto $X_4$;
    \item Defina $S_p = X_4 - X_1$;
    \item Faça uma busca unidimensional ao longo de $S_p$;
    \item Mostre (em quadro) que, neste exemplo quadrático, o ponto ótimo é atingido rapidamente.
\end{itemize}

\textbf{Mensagem didática:}
\begin{itemize}
    \item Powell é um método direto eficiente para funções suaves;
    \item Para funções quadráticas, é teoricamente quadraticamente convergente;
    \item Em funções gerais, continua sendo útil, embora sem as garantias exatas de convergência em $n$ ciclos.
\end{itemize}

\section{Método do Simplex (Nelder--Mead)}

\subsection{Visão Geométrica}

Definição de simplex:
\begin{itemize}
    \item Em $\mathbb{R}^n$, um \textbf{simplex} é o envelope convexo de $n+1$ pontos não coplanares;
    \item Para $n=1$, é um segmento de reta (2 pontos);
    \item Para $n=2$, é um triângulo (3 pontos);
    \item Para $n=3$, é um tetraedro (4 pontos), etc.
\end{itemize}

No método de Nelder--Mead:
\begin{itemize}
    \item Mantém-se um simplex a cada iteração;
    \item Em cada passo, avalia-se $f$ nos vértices;
    \item Movimenta-se o simplex via:
    \begin{itemize}
        \item \textbf{Reflexão} do pior vértice;
        \item \textbf{Expansão}, se a reflexão produziu um valor muito bom;
        \item \textbf{Contração}, se a reflexão não foi boa;
        \item \textbf{Encolhimento} (\emph{shrink}), se mesmo a contração falhar.
    \end{itemize}
\end{itemize}

\subsection{Passos Básicos em 2D (Triângulo)}

Considere $n=2$ (didaticamente mais simples):

\begin{enumerate}
    \item Iniciar com três pontos $X_1, X_2, X_3$ e calcular $f(X_i)$.
    \item Identificar:
    \begin{itemize}
        \item $X_l$: ponto com menor valor $f$ (melhor);
        \item $X_h$: ponto com maior valor $f$ (pior);
        \item $X_m$: ponto intermediário.
    \end{itemize}
    \item Calcular o centroide dos vértices exceto o pior:
    \[
        X_0 = \frac{1}{n} \sum_{i \neq h} X_i.
    \]
    \item \textbf{Reflexão}:
    \[
        X_r = X_0 + \alpha (X_0 - X_h),\quad \alpha>0.
    \]
    \item Avaliar $f(X_r)$ e aplicar a regra:
    \begin{itemize}
        \item Se $f(X_l) \le f(X_r) < f(X_m)$, substituir $X_h$ por $X_r$ e seguir para a próxima iteração;
        \item Se $f(X_r) < f(X_l)$, tentar \textbf{expansão};
        \item Se $f(X_r) \ge f(X_m)$, tentar \textbf{contração}.
    \end{itemize}
    \item \textbf{Expansão} (se $X_r$ for muito bom):
    \[
        X_e = X_0 + \gamma (X_r - X_0),\quad \gamma > 1.
    \]
    Escolher entre $X_r$ e $X_e$ o ponto com menor valor para substituir $X_h$.
    \item \textbf{Contração}: se reflexão foi ruim:
    \[
        X_c = X_0 + \beta (X_h - X_0),\quad 0<\beta<1.
    \]
    Se $f(X_c)$ melhorar o simplex, substituir $X_h$ por $X_c$. Caso contrário, realizar \textbf{encolhimento}:
    \[
        X_i \leftarrow \frac{X_i + X_l}{2},\quad \text{para todos os vértices } i.
    \]
\end{enumerate}

\subsection{Exemplo Numérico Esquemático}

Considere novamente
\[
    f(x_1,x_2) = x_1 - x_2 + 2x_1^2 + 2x_1x_2 + x_2^2,
\]
e um simplex inicial com:
\[
    X_1 = (4,4),\quad
    X_2 = (5,4),\quad
    X_3 = (4,5).
\]

Sugestão para a aula:
\begin{enumerate}
    \item Calcular $f(X_1)$, $f(X_2)$, $f(X_3)$;
    \item Identificar $X_l$, $X_m$, $X_h$;
    \item Calcular $X_0$ e construir $X_r$;
    \item Verificar se é caso de expansão ou não;
    \item Desenhar o triângulo original e o triângulo após a operação em um gráfico de contorno (mesmo qualitativo).
\end{enumerate}

\subsection{Comentários Finais sobre o Simplex de Nelder--Mead}

\begin{itemize}
    \item Método direto popular em baixa dimensão;
    \item Não usa derivadas, é puramente geométrico;
    \item Funciona bem em muitas aplicações práticas, mas não tem garantias fortes de convergência global;
    \item Em dimensões mais altas, o comportamento pode se deteriorar (necessidade de muitas avaliações de $f$ e geometria do simplex pode ficar pobre).
\end{itemize}

\section*{Resumo e Pontos para Fixação}

\begin{itemize}
    \item \textbf{Métodos diretos} são apropriados quando derivadas são indisponíveis, não confiáveis ou muito caras.
    \item Métodos simples como \textbf{Random Jumping}, \textbf{Random Walk} e \textbf{Grid Search} são úteis para explorar o espaço e encontrar regiões promissoras.
    \item O \textbf{método univariado} é conceitualmente simples, mas pode ser lento e sofrer com vales estreitos.
    \item O \textbf{método de Powell} explora direções padrão (aproximadamente conjugadas) para acelerar a convergência em funções suaves, especialmente quadráticas.
    \item O \textbf{método do simplex (Nelder--Mead)} usa operações geométricas (reflexão, expansão, contração, encolhimento) sobre um simplex para aproximar o mínimo.
\end{itemize}

\section*{Sugestões de Exercícios}

\begin{enumerate}
    \item Implementar o método univariado para a função
    \[
        f(x_1,x_2) = x_1 - x_2 + 2x_1^2 + 2x_1x_2 + x_2^2
    \]
    e comparar a trajetória com um método de gradiente.
    \item Implementar o método de Powell em duas dimensões e testar em:
    \begin{itemize}
        \item uma função quadrática;
        \item uma função não quadrática com vale alongado.
    \end{itemize}
    \item Implementar o método do simplex de Nelder--Mead em 2D, reproduzindo pelo menos duas iterações completas (com reflexão e alguma combinação de expansão/contração).
    \item Investigar numericamente o efeito da escolha dos parâmetros $\alpha$, $\beta$ e $\gamma$ no comportamento do método do simplex.
\end{enumerate}

\end{document}
